% Target ~120 words. Replaces 25-word stub.

We presented Articulatory Trajectory Deviation (ATD), a frozen-SSL
articulatory metric with within-corpus calibration for dysarthric
speech severity that exploits the geometric structure of frozen
self-supervised speech representations. ATD attains speaker-level
Spearman $\rho{=}0.756$ on TORGO at WavLM Layer~9 without fine-tuning ---
a figure that is unchanged under leave-one-control-out leakage control,
where the early-layer alternative inflates. A fairly-tuned linear probe on
mean-pooled frozen-WavLM features, refit per leave-one-speaker-out fold,
matches this accuracy ($\rho{=}0.760$--$0.775$) and attains a lower RMSE
($0.581$ vs.\ $0.671$); ATD is therefore \emph{not} more accurate than a
trained probe, and its contribution is comparable accuracy \emph{without
any training or labels}, together with per-phoneme interpretability and
real-time operation. ATD does beat a fairly-tuned MFCC baseline
($\rho{=}0.609$). A pre-registered controlled-
perturbation study \emph{falsified} our articulatory-specificity
hypothesis (additive noise perturbs ATD more than formant shift does),
which we report rather than bury; the study's surviving positive is that
the trajectory metric detects frame dropout better than pooled MFCC.
Phoneme-level DTW path decomposition (into per-phoneme contributions plus
a silence/unattributed residual that carries the majority of the path
cost on single-word audio) extends the headline scalar to per-token
deviation. The frozen SSL articulatory manifold is shown to encode a
measurable, falsifiable, and reproducible severity signal without
requiring labelled training data. Genuine zero-shot cross-corpus transfer
failed (TORGO$\to$IPVS $\rho{=}0.18$, NS; TORGO$\to$UASpeech
$\rho{=}-0.22$), so we position the method as within-corpus calibration,
not a zero-shot metric. Critically,
a pooled 4-corpus analysis on 28{,}082 utterances from 164 speakers
shows ATD is specific to dysarthric pathology rather than pronunciation
deviation in general: across 24 L2-ARCTIC speakers the correlation
between mean ATD and ground-truth phone-error rate is null
($\rho{=}0.041$, $p{=}0.85$), and across 110 VCTK speakers spanning 13
native English accents Kruskal--Wallis is not significant
($p{=}0.075$), while the dysarthric-vs-healthy contrast on the pooled
corpus remains highly significant ($\hat{\beta}{=}0.75$ z-units,
$p{=}8.9{\times}10^{-5}$). ATD measures pathology, not accent.
